\section{家庭、劳作与知识的万历证据}

本组材料把视线从朝廷决策移向家庭、劳作、环境与知识如何在文书中留下痕迹。土地、户籍、税役、河工、学校、宗教空间与地方灾情，往往只在征解、工程、诉讼、捐输或褒奖的片段中出现。材料的缺口并非无关紧要：它说明哪些人的行动更容易进入档案，哪些人的经验只能从地方记录和物质遗存中间接追索。

\subsection{清丈、纳银与地方家庭}

\eventanchor{ML-1577-001}{明·1577（万历五年）}
\paragraph{1577（万历五年）：辽东、西北、朝鲜或西南军务的本年调兵与战报记录跨区域动员，补给与兵额需要独立核验} 这一锚点使家庭、劳动或地方资源进入可追查的历史时间。账本注明的把握范围是来源导向的年度桥接条目：本年材料可定位相关政务线索；具体月日、发文机关、地域和执行结果仍须逐项回查，不能以制度文本或奏报替代实际效果。《神宗实录》万历五年条；《明史·兵志》及相关列传；边镇、地方志或对方材料 提供了定位事件的材料链，却分别反映编纂者、报送者和保存者的视角。其发生条件可概括为前期边防、地方武装或跨区域资源调动的压力推动了该行动；随后可见的制度或社会影响是它改变了后续调兵、补给或地方秩序的条件；战果和伤亡须分开核对。对于日常生活、性别与家庭分工、非精英劳作或未留名者，仍不能由这些材料直接补足。桥接条目只标示可回查的年度问题与材料链；实录有中央选择性，崇祯期尤其需要奏疏、地方、朝鲜及满文材料交叉。

\eventanchor{ML-1577-002}{明·1577（万历五年）}
\paragraph{1577（万历五年）：朝鲜战争、海上贸易和朝贡使节的本年多方记录提供跨政权比较，外交语言不能替代地方经验} 这一锚点使家庭、劳动或地方资源进入可追查的历史时间。账本注明的把握范围是来源导向的年度桥接条目：本年材料可定位相关政务线索；具体月日、发文机关、地域和执行结果仍须逐项回查，不能以制度文本或奏报替代实际效果。《神宗实录》万历五年条；《明史·外国传》；《朝鲜王朝实录》、琉球《历代宝案》或海外档案 提供了定位事件的材料链，却分别反映编纂者、报送者和保存者的视角。其发生条件可概括为朝贡名义、边境安全与港口贸易的利益使交涉持续发生；随后可见的制度或社会影响是它为后续册封、互市或海防安排留下文书与跨政权比较线索。对于日常生活、性别与家庭分工、非精英劳作或未留名者，仍不能由这些材料直接补足。桥接条目只标示可回查的年度问题与材料链；实录有中央选择性，崇祯期尤其需要奏疏、地方、朝鲜及满文材料交叉。

\eventanchor{ML-1577-003}{明·1577（万历五年）}
\paragraph{1577（万历五年）：清丈、一条鞭、漕运或军饷的本年文书构成万历财政的可核查节点，定额、报数和实收必须分列} 这一锚点使家庭、劳动或地方资源进入可追查的历史时间。账本注明的把握范围是来源导向的年度桥接条目：本年材料可定位相关政务线索；具体月日、发文机关、地域和执行结果仍须逐项回查，不能以制度文本或奏报替代实际效果。《神宗实录》万历五年条；《明会典》相关条目；地方志、黄册/契约/河工材料 提供了定位事件的材料链，却分别反映编纂者、报送者和保存者的视角。其发生条件可概括为财政征解、交通工程或货币支付的压力使相关措施进入政务；随后可见的制度或社会影响是其法定安排不等于地方执行，后续须以地域材料追踪负担与成效。对于日常生活、性别与家庭分工、非精英劳作或未留名者，仍不能由这些材料直接补足。桥接条目只标示可回查的年度问题与材料链；实录有中央选择性，崇祯期尤其需要奏疏、地方、朝鲜及满文材料交叉。

\eventanchor{ML-1577-004}{明·1577（万历五年）}
\paragraph{1577（万历五年）：书院、科举、出版与宗教活动的本年线索可连接知识网络，存世文本有阶层与地域偏差} 这一锚点使家庭、劳动或地方资源进入可追查的历史时间。账本注明的把握范围是来源导向的年度桥接条目：本年材料可定位相关政务线索；具体月日、发文机关、地域和执行结果仍须逐项回查，不能以制度文本或奏报替代实际效果。《神宗实录》万历五年条；《明史·选举志》或《艺文志》；文集、碑刻或书目材料 提供了定位事件的材料链，却分别反映编纂者、报送者和保存者的视角。其发生条件可概括为制度、教育或知识传播的需要推动了相关行动或文本形成；随后可见的制度或社会影响是它提供制度和传播史的锚点，不能据此推断社会整体接受。对于日常生活、性别与家庭分工、非精英劳作或未留名者，仍不能由这些材料直接补足。桥接条目只标示可回查的年度问题与材料链；实录有中央选择性，崇祯期尤其需要奏疏、地方、朝鲜及满文材料交叉。

\eventanchor{MM-1577-PAN-JIXUN-RIVER-CONTROL}{明·1577（万历五年）}
\paragraph{1577（万历五年）：潘季驯受命主持河防，提出束水攻沙并组织黄河堤防与运河保运工程} 这一锚点使家庭、劳动或地方资源进入可追查的历史时间。账本注明的把握范围是任命、奏议和工程命令可靠；治河成效随河段、水情和维护周期而异，后出技术赞誉不能证明所有堤段安全或漕运畅通。《神宗实录》万历五年河道条；潘季驯《河防一览》；《明史·河渠志》 提供了定位事件的材料链，却分别反映编纂者、报送者和保存者的视角。其发生条件可概括为黄河决溢威胁运河漕运、沿河农田和漕粮北运，中央需在河工经费、役夫征发与水势变化之间选择方案；随后可见的制度或社会影响是堤防与导流工程加强了对若干河段的控制并保障部分漕运，亦持续要求地方承担经费和劳役；河患不因一次治理而消除。对于日常生活、性别与家庭分工、非精英劳作或未留名者，仍不能由这些材料直接补足。因此，登记、报送与实际经验之间仍须保留距离。

\eventanchor{MM-1577-ZHANG-JUZHENG-DUOQING}{明·1577（万历五年）}
\paragraph{1577（万历五年）：张居正父丧期间奉旨夺情留任，引发朝臣围绕礼制与政务的争论} 这一锚点使家庭、劳动或地方资源进入可追查的历史时间。账本注明的把握范围是夺情诏令和争论可靠；反对者并非只出于礼制或派系，具体立场应回到各人奏疏。《神宗实录》万历五年条；《明史·张居正传》；张居正、吴中行等奏疏 提供了定位事件的材料链，却分别反映编纂者、报送者和保存者的视角。其发生条件可概括为张居正主持的改革需要连续推动，太后与皇帝担忧其离职；传统士大夫伦理则强调居丧服制；随后可见的制度或社会影响是张居正得以留任，但政治反对扩大；礼制争论成为评价其权力边界和万历初政的重要线索。对于日常生活、性别与家庭分工、非精英劳作或未留名者，仍不能由这些材料直接补足。因此，登记、报送与实际经验之间仍须保留距离。

\subsection{河工、军役与劳作}

\eventanchor{ML-1578-001}{明·1578（万历六年）}
\paragraph{1578（万历六年）：葡萄牙人获准前往广州} 这一锚点使家庭、劳动或地方资源进入可追查的历史时间。账本注明的把握范围是编年候选的年序可由官修与后出材料交叉；具体月日、参与者、战果或数字必须回查所列材料，不能将单一叙事当作定论。《神宗实录》万历六年条；《明史·外国传》；《朝鲜王朝实录》、琉球《历代宝案》或海外档案 提供了定位事件的材料链，却分别反映编纂者、报送者和保存者的视角。其发生条件可概括为朝贡名义、边境安全与港口贸易的利益使交涉持续发生；随后可见的制度或社会影响是它为后续册封、互市或海防安排留下文书与跨政权比较线索。对于日常生活、性别与家庭分工、非精英劳作或未留名者，仍不能由这些材料直接补足。译写候选以编年材料定位；实录记载反映朝廷选择，崇祯期须另以奏疏、地方及敌对政权材料交叉。

\eventanchor{ML-1578-002}{明·1578（万历六年）}
\paragraph{1578（万历六年）：书院、科举、出版与宗教活动的本年线索可连接知识网络，存世文本有阶层与地域偏差} 这一锚点使家庭、劳动或地方资源进入可追查的历史时间。账本注明的把握范围是来源导向的年度桥接条目：本年材料可定位相关政务线索；具体月日、发文机关、地域和执行结果仍须逐项回查，不能以制度文本或奏报替代实际效果。《神宗实录》万历六年条；《明史·选举志》或《艺文志》；文集、碑刻或书目材料 提供了定位事件的材料链，却分别反映编纂者、报送者和保存者的视角。其发生条件可概括为制度、教育或知识传播的需要推动了相关行动或文本形成；随后可见的制度或社会影响是它提供制度和传播史的锚点，不能据此推断社会整体接受。对于日常生活、性别与家庭分工、非精英劳作或未留名者，仍不能由这些材料直接补足。桥接条目只标示可回查的年度问题与材料链；实录有中央选择性，崇祯期尤其需要奏疏、地方、朝鲜及满文材料交叉。

\eventanchor{ML-1578-003}{明·1578（万历六年）}
\paragraph{1578（万历六年）：考成、人事、立储和留中等本年材料标记皇帝、内阁与言官的协调关系，不可只用党争标签解释} 这一锚点使家庭、劳动或地方资源进入可追查的历史时间。账本注明的把握范围是来源导向的年度桥接条目：本年材料可定位相关政务线索；具体月日、发文机关、地域和执行结果仍须逐项回查，不能以制度文本或奏报替代实际效果。《神宗实录》万历六年条；《明史·本纪》及相关列传；诏令、奏疏或会典版本 提供了定位事件的材料链，却分别反映编纂者、报送者和保存者的视角。其发生条件可概括为继承、官僚协调或皇权运作的压力使问题进入朝廷程序；随后可见的制度或社会影响是它影响后续人事与政策议程；动机和派系标签不能由单一叙事裁定。对于日常生活、性别与家庭分工、非精英劳作或未留名者，仍不能由这些材料直接补足。桥接条目只标示可回查的年度问题与材料链；实录有中央选择性，崇祯期尤其需要奏疏、地方、朝鲜及满文材料交叉。

\eventanchor{ML-1578-004}{明·1578（万历六年）}
\paragraph{1578（万历六年）：本年灾荒、河工和地方赈济材料可用于研究风险分配，灾害叙事和统计数字应区分} 这一锚点使家庭、劳动或地方资源进入可追查的历史时间。账本注明的把握范围是来源导向的年度桥接条目：本年材料可定位相关政务线索；具体月日、发文机关、地域和执行结果仍须逐项回查，不能以制度文本或奏报替代实际效果。《神宗实录》万历六年条；《明史·五行志》；受灾府县地方志与奏疏 提供了定位事件的材料链，却分别反映编纂者、报送者和保存者的视角。其发生条件可概括为自然风险与地方报告促成朝廷或地方的应对记录；随后可见的制度或社会影响是赈济、迁徙和损失的实际范围仍须以受灾地区材料分别核验。对于日常生活、性别与家庭分工、非精英劳作或未留名者，仍不能由这些材料直接补足。桥接条目只标示可回查的年度问题与材料链；实录有中央选择性，崇祯期尤其需要奏疏、地方、朝鲜及满文材料交叉。

\paragraph{1578（万历六年）六月：万历帝册立王氏为皇后，完成幼帝成年后的皇后册立仪式} 这一锚点使家庭、劳动或地方资源进入可追查的历史时间。账本注明的把握范围是册立日期、礼仪和诏令可靠；礼制记录突出朝廷秩序，不能据此推知帝后关系、后宫实际权力或后继生育结果。《神宗实录》万历六年六月乙卯条；《明史·孝端显皇后传》；明代后妃册立仪文 提供了定位事件的材料链，却分别反映编纂者、报送者和保存者的视角。其发生条件可概括为万历帝已亲政，皇室需要依礼确定皇后以组织后宫名分、祭祀和未来继嗣的程序；随后可见的制度或社会影响是王氏取得皇后名号，后宫礼制与外戚关系有了正式支点；继承问题仍取决于子嗣、册立和宫廷政治。对于日常生活、性别与家庭分工、非精英劳作或未留名者，仍不能由这些材料直接补足。因此，登记、报送与实际经验之间仍须保留距离。

\paragraph{1578（万历六年）：朝廷令各省清丈田亩、核查隐漏，为赋役重整收集册籍} 这一锚点使家庭、劳动或地方资源进入可追查的历史时间。账本注明的把握范围是清丈命令可靠；各省启动、完结与报亩年份不同，中央汇总数不能当作同一年度的全国实测。《神宗实录》万历六年户部条；《明史·食货志》；徽州、江南等地清丈册与契约 提供了定位事件的材料链，却分别反映编纂者、报送者和保存者的视角。其发生条件可概括为旧有登记与实际占有、赋额脱节，中央财政和改革政策需要更可比的田亩、户役信息；随后可见的制度或社会影响是各地重编册籍并暴露隐漏、产权和役额争议；清丈为一条鞭的区域推行提供条件，却不决定其统一形式。对于日常生活、性别与家庭分工、非精英劳作或未留名者，仍不能由这些材料直接补足。因此，登记、报送与实际经验之间仍须保留距离。

\subsection{灾荒、迁徙与地方互助}

\paragraph{1579（万历七年）：东林运动：私立东林书院全部关闭} 这一锚点使家庭、劳动或地方资源进入可追查的历史时间。账本注明的把握范围是编年候选的年序可由官修与后出材料交叉；具体月日、参与者、战果或数字必须回查所列材料，不能将单一叙事当作定论。《神宗实录》万历七年条；《明史·本纪》及相关列传；诏令、奏疏或会典版本 提供了定位事件的材料链，却分别反映编纂者、报送者和保存者的视角。其发生条件可概括为继承、官僚协调或皇权运作的压力使问题进入朝廷程序；随后可见的制度或社会影响是它影响后续人事与政策议程；动机和派系标签不能由单一叙事裁定。对于日常生活、性别与家庭分工、非精英劳作或未留名者，仍不能由这些材料直接补足。译写候选以编年材料定位；实录记载反映朝廷选择，崇祯期须另以奏疏、地方及敌对政权材料交叉。

\paragraph{1579（万历七年）：朝鲜战争、海上贸易和朝贡使节的本年多方记录提供跨政权比较，外交语言不能替代地方经验} 这一锚点使家庭、劳动或地方资源进入可追查的历史时间。账本注明的把握范围是来源导向的年度桥接条目：本年材料可定位相关政务线索；具体月日、发文机关、地域和执行结果仍须逐项回查，不能以制度文本或奏报替代实际效果。《神宗实录》万历七年条；《明史·外国传》；《朝鲜王朝实录》、琉球《历代宝案》或海外档案 提供了定位事件的材料链，却分别反映编纂者、报送者和保存者的视角。其发生条件可概括为朝贡名义、边境安全与港口贸易的利益使交涉持续发生；随后可见的制度或社会影响是它为后续册封、互市或海防安排留下文书与跨政权比较线索。对于日常生活、性别与家庭分工、非精英劳作或未留名者，仍不能由这些材料直接补足。桥接条目只标示可回查的年度问题与材料链；实录有中央选择性，崇祯期尤其需要奏疏、地方、朝鲜及满文材料交叉。

\paragraph{1579（万历七年）：本年灾荒、河工和地方赈济材料可用于研究风险分配，灾害叙事和统计数字应区分} 这一锚点使家庭、劳动或地方资源进入可追查的历史时间。账本注明的把握范围是来源导向的年度桥接条目：本年材料可定位相关政务线索；具体月日、发文机关、地域和执行结果仍须逐项回查，不能以制度文本或奏报替代实际效果。《神宗实录》万历七年条；《明史·五行志》；受灾府县地方志与奏疏 提供了定位事件的材料链，却分别反映编纂者、报送者和保存者的视角。其发生条件可概括为自然风险与地方报告促成朝廷或地方的应对记录；随后可见的制度或社会影响是赈济、迁徙和损失的实际范围仍须以受灾地区材料分别核验。对于日常生活、性别与家庭分工、非精英劳作或未留名者，仍不能由这些材料直接补足。桥接条目只标示可回查的年度问题与材料链；实录有中央选择性，崇祯期尤其需要奏疏、地方、朝鲜及满文材料交叉。

\paragraph{1579（万历七年）：辽东、西北、朝鲜或西南军务的本年调兵与战报记录跨区域动员，补给与兵额需要独立核验} 这一锚点使家庭、劳动或地方资源进入可追查的历史时间。账本注明的把握范围是来源导向的年度桥接条目：本年材料可定位相关政务线索；具体月日、发文机关、地域和执行结果仍须逐项回查，不能以制度文本或奏报替代实际效果。《神宗实录》万历七年条；《明史·兵志》及相关列传；边镇、地方志或对方材料 提供了定位事件的材料链，却分别反映编纂者、报送者和保存者的视角。其发生条件可概括为前期边防、地方武装或跨区域资源调动的压力推动了该行动；随后可见的制度或社会影响是它改变了后续调兵、补给或地方秩序的条件；战果和伤亡须分开核对。对于日常生活、性别与家庭分工、非精英劳作或未留名者，仍不能由这些材料直接补足。桥接条目只标示可回查的年度问题与材料链；实录有中央选择性，崇祯期尤其需要奏疏、地方、朝鲜及满文材料交叉。

\paragraph{1579（万历七年）：清丈、一条鞭、漕运或军饷的本年文书构成万历财政的可核查节点，定额、报数和实收必须分列} 这一锚点使家庭、劳动或地方资源进入可追查的历史时间。账本注明的把握范围是来源导向的年度桥接条目：本年材料可定位相关政务线索；具体月日、发文机关、地域和执行结果仍须逐项回查，不能以制度文本或奏报替代实际效果。《神宗实录》万历七年条；《明会典》相关条目；地方志、黄册/契约/河工材料 提供了定位事件的材料链，却分别反映编纂者、报送者和保存者的视角。其发生条件可概括为财政征解、交通工程或货币支付的压力使相关措施进入政务；随后可见的制度或社会影响是其法定安排不等于地方执行，后续须以地域材料追踪负担与成效。对于日常生活、性别与家庭分工、非精英劳作或未留名者，仍不能由这些材料直接补足。桥接条目只标示可回查的年度问题与材料链；实录有中央选择性，崇祯期尤其需要奏疏、地方、朝鲜及满文材料交叉。

\paragraph{1579（万历七年）：张居正政府以禁讲学、限制书院活动的诏令约束地方士人结社} 这一锚点使家庭、劳动或地方资源进入可追查的历史时间。账本注明的把握范围是禁令和中央政策方向可靠；各地书院是否停讲、改名或私下延续须以地方志、碑刻和文集逐处核验，不能以诏令写成思想生活完全中断。《神宗实录》万历七年书院条；张居正公牍与奏疏；书院碑记及地方志 提供了定位事件的材料链，却分别反映编纂者、报送者和保存者的视角。其发生条件可概括为张居正担忧讲学社群与地方结交影响官员考成和政治批评，中央希望将士人活动重新纳入官学与官僚秩序；随后可见的制度或社会影响是部分书院公开活动受限，士人转向家塾、文会或其他交往形式；禁令也为后来的学术自治记忆提供素材。对于日常生活、性别与家庭分工、非精英劳作或未留名者，仍不能由这些材料直接补足。因此，登记、报送与实际经验之间仍须保留距离。

\subsection{书院、刻书与知识网络}

\paragraph{1580（万历八年）：单鞭法：简化税法} 这一锚点使家庭、劳动或地方资源进入可追查的历史时间。账本注明的把握范围是编年候选的年序可由官修与后出材料交叉；具体月日、参与者、战果或数字必须回查所列材料，不能将单一叙事当作定论。《神宗实录》万历八年条；《明会典》相关条目；地方志、黄册/契约/河工材料 提供了定位事件的材料链，却分别反映编纂者、报送者和保存者的视角。其发生条件可概括为财政征解、交通工程或货币支付的压力使相关措施进入政务；随后可见的制度或社会影响是其法定安排不等于地方执行，后续须以地域材料追踪负担与成效。对于日常生活、性别与家庭分工、非精英劳作或未留名者，仍不能由这些材料直接补足。译写候选以编年材料定位；实录记载反映朝廷选择，崇祯期须另以奏疏、地方及敌对政权材料交叉。

\paragraph{1580（万历八年）：官员们批评万历皇帝的疏忽和个人生活的不当行为} 这一锚点使家庭、劳动或地方资源进入可追查的历史时间。账本注明的把握范围是编年候选的年序可由官修与后出材料交叉；具体月日、参与者、战果或数字必须回查所列材料，不能将单一叙事当作定论。《神宗实录》万历八年条；《明史·本纪》及相关列传；诏令、奏疏或会典版本 提供了定位事件的材料链，却分别反映编纂者、报送者和保存者的视角。其发生条件可概括为继承、官僚协调或皇权运作的压力使问题进入朝廷程序；随后可见的制度或社会影响是它影响后续人事与政策议程；动机和派系标签不能由单一叙事裁定。对于日常生活、性别与家庭分工、非精英劳作或未留名者，仍不能由这些材料直接补足。译写候选以编年材料定位；实录记载反映朝廷选择，崇祯期须另以奏疏、地方及敌对政权材料交叉。

\paragraph{1580（万历八年）：清丈、一条鞭、漕运或军饷的本年文书构成万历财政的可核查节点，定额、报数和实收必须分列} 这一锚点使家庭、劳动或地方资源进入可追查的历史时间。账本注明的把握范围是来源导向的年度桥接条目：本年材料可定位相关政务线索；具体月日、发文机关、地域和执行结果仍须逐项回查，不能以制度文本或奏报替代实际效果。《神宗实录》万历八年条；《明会典》相关条目；地方志、黄册/契约/河工材料 提供了定位事件的材料链，却分别反映编纂者、报送者和保存者的视角。其发生条件可概括为财政征解、交通工程或货币支付的压力使相关措施进入政务；随后可见的制度或社会影响是其法定安排不等于地方执行，后续须以地域材料追踪负担与成效。对于日常生活、性别与家庭分工、非精英劳作或未留名者，仍不能由这些材料直接补足。桥接条目只标示可回查的年度问题与材料链；实录有中央选择性，崇祯期尤其需要奏疏、地方、朝鲜及满文材料交叉。

\paragraph{1580（万历八年）：朝鲜战争、海上贸易和朝贡使节的本年多方记录提供跨政权比较，外交语言不能替代地方经验} 这一锚点使家庭、劳动或地方资源进入可追查的历史时间。账本注明的把握范围是来源导向的年度桥接条目：本年材料可定位相关政务线索；具体月日、发文机关、地域和执行结果仍须逐项回查，不能以制度文本或奏报替代实际效果。《神宗实录》万历八年条；《明史·外国传》；《朝鲜王朝实录》、琉球《历代宝案》或海外档案 提供了定位事件的材料链，却分别反映编纂者、报送者和保存者的视角。其发生条件可概括为朝贡名义、边境安全与港口贸易的利益使交涉持续发生；随后可见的制度或社会影响是它为后续册封、互市或海防安排留下文书与跨政权比较线索。对于日常生活、性别与家庭分工、非精英劳作或未留名者，仍不能由这些材料直接补足。桥接条目只标示可回查的年度问题与材料链；实录有中央选择性，崇祯期尤其需要奏疏、地方、朝鲜及满文材料交叉。

\paragraph{1580（万历八年）：考成、人事、立储和留中等本年材料标记皇帝、内阁与言官的协调关系，不可只用党争标签解释} 这一锚点使家庭、劳动或地方资源进入可追查的历史时间。账本注明的把握范围是来源导向的年度桥接条目：本年材料可定位相关政务线索；具体月日、发文机关、地域和执行结果仍须逐项回查，不能以制度文本或奏报替代实际效果。《神宗实录》万历八年条；《明史·本纪》及相关列传；诏令、奏疏或会典版本 提供了定位事件的材料链，却分别反映编纂者、报送者和保存者的视角。其发生条件可概括为继承、官僚协调或皇权运作的压力使问题进入朝廷程序；随后可见的制度或社会影响是它影响后续人事与政策议程；动机和派系标签不能由单一叙事裁定。对于日常生活、性别与家庭分工、非精英劳作或未留名者，仍不能由这些材料直接补足。桥接条目只标示可回查的年度问题与材料链；实录有中央选择性，崇祯期尤其需要奏疏、地方、朝鲜及满文材料交叉。

\paragraph{1580（万历八年）：朝廷以清丈和赋役整顿为基础，推动各地合并田赋、丁役与杂派的征收办法} 这一锚点使家庭、劳动或地方资源进入可追查的历史时间。账本注明的把握范围是中央推动赋役合并的方向可靠；“一条鞭法”在各省府县的项目、折银比例和实施节奏并不一致。《神宗实录》万历八年户部条；《明史·食货志》；各地赋役全书、清丈册和契约 提供了定位事件的材料链，却分别反映编纂者、报送者和保存者的视角。其发生条件可概括为清丈揭示登记与赋役的脱节，役法名目繁复且征解成本高，张居正政府希望简化责任与提高可核算性；随后可见的制度或社会影响是更多地区以田亩为主要计征基础并折纳部分银两；不同地方仍保留实物、差役或本地变通。对于日常生活、性别与家庭分工、非精英劳作或未留名者，仍不能由这些材料直接补足。因此，登记、报送与实际经验之间仍须保留距离。
